\documentclass[11pt,a4paper]{article}
\usepackage[T1]{fontenc}
\usepackage{lmodern}
\usepackage[margin=27mm]{geometry}
\usepackage{amsmath,amssymb,amsthm}
\usepackage{microtype}
\usepackage{enumitem}
\usepackage[hidelinks]{hyperref}
\hypersetup{pdftitle={The first missing multiple in a finite set},pdfsubject={Matching asymptotic bounds for the sum of first missing multipliers},pdfkeywords={number theory, extremal combinatorics, compression}}
\setlist{itemsep=2pt,topsep=5pt}
\setlength{\emergencystretch}{2em}
\newtheorem{theorem}{Theorem}
\newtheorem{lemma}{Lemma}
\newcommand{\N}{\mathbb N}
\newcommand{\Nzero}{\mathbb N_0}
\newcommand{\rad}{\operatorname{rad}}

\begin{document}
\begin{center}
{\Large The first missing multiple in a finite set}
\end{center}
\vspace{0.5em}
\begin{abstract}
For a finite set $S$ of positive integers, let $k_S(a)$ be the least positive integer $k$ for which $ka\notin S$.
We prove that the maximum of $\sum_{a\in S}k_S(a)$ over sets of cardinality $n$ has order $n\log n\log\log n$.
The upper bound follows by compressing prime exponents and double counting prime-multiplication pairs.
A construction with separate constraints on small and large prime factors gives the matching lower bound.
\end{abstract}

\begin{theorem}\label{thm:main}
There are absolute constants $c,C>0$ such that, for all sufficiently large $n$,
\[
c\,n\log n\log\log n
\ \le\ 
\max_{\substack{S\subset\N\\|S|=n}}\ \sum_{a\in S}k_S(a)
\ \le\ 
C\,n\log n\log\log n.
\]
\end{theorem}

All logarithms are natural unless a base is indicated.
The set $S=\{1,\ldots,n\}$ gives
\[
\sum_{a\in S}k_S(a)
=n+\sum_{a=1}^n\left\lfloor\frac na\right\rfloor
=n\log n+O(n),
\]
so it does not attain the correct order of growth.

Write $r_S(a)=k_S(a)-1$ for $a\in S$.
For $t\ge1$, define
\[
N_t(S)=\bigl|\{a\in\N:a,2a,\ldots,ta\in S\}\bigr|.
\]
Then
\begin{equation}\label{eq:tails}
\sum_{a\in S}k_S(a)=|S|+\sum_{t\ge1}N_t(S).
\end{equation}
The sums are finite because $r_S(a)\le |S|$.
We use the classical Chebyshev estimates
\begin{equation}\label{eq:primes}
\vartheta(x):=\sum_{p\le x}\log p\asymp x,
\qquad
\pi(x)\asymp\frac{x}{\log x},
\end{equation}
where sums over $p$ run over primes.
In particular, for the $j$-th prime $p_j$,
\begin{equation}\label{eq:nth-prime}
\log p_j=O(\log(2j)).
\end{equation}
An elementary justification of these estimates is included at the end.

\section{Compression}

\begin{lemma}\label{lem:intervals}
Let $E_1,\ldots,E_s$ be finite subsets of $\Nzero$ and let $L_1,\ldots,L_s$ be nonnegative integers.
Replacing each $E_i$ by $\{0,\ldots,|E_i|-1\}$ cannot decrease the number of integers $e\ge0$ satisfying
\[
e+\{0,\ldots,L_i\}\subseteq E_i\qquad(1\le i\le s).
\]
An initial interval of length zero is understood to be empty.
\end{lemma}

\begin{proof}
For one set $E$ and one length $L$, there are at most $\max(|E|-L,0)$ valid starting positions.
Indeed, if there are any, all valid starts together with the $L$ integers following the largest valid start are distinct elements of $E$.
Thus the number satisfying all requirements is at most
\[
\min_i\max(|E_i|-L_i,0).
\]
After replacement, the valid starts for the $i$-th requirement form the initial interval of exactly this individual length.
Their intersection attains the displayed upper bound.
\end{proof}

\newpage
\begin{lemma}\label{lem:compression}
Every finite $S\subset\N$ can be replaced by a set $T$ of the same cardinality such that $N_t(T)\ge N_t(S)$ for every $t\ge1$, and:
\begin{enumerate}[label=(\roman*)]
\item every divisor of an element of $T$ belongs to $T$;
\item if $a\in T$, $q\mid a$, and $p<q$ are primes, then $ap/q\in T$.
\end{enumerate}
\end{lemma}

\begin{proof}
We use two operations, each applied simultaneously to all chains in the indicated partition of $\N$.

For a prime $p$, consider the chains
\[
b,bp,bp^2,\ldots,\qquad p\nmid b.
\]
If a chain contains $s$ elements of the current set, replace them by its first $s$ elements.
Fix $t$ and a source chain, and write a possible source element as $a=bp^e$.
For each $c\le t$ with $p\nmid c$, the multipliers $cp^h\le t$ have exponents $h=0,\ldots,L_c$.
The requirements $a(cp^h)\in S$ therefore say
\[
e+\{0,\ldots,L_c\}\subseteq E_c,
\qquad
E_c=\{u\ge0:bc\,p^u\in S\}.
\]
Since $p\nmid bc$, each $E_c$ is the occupied set of a chain in this partition, and distinct $c$ give distinct target chains.
Compression replaces every $E_c$ by an initial interval of the same size.
Lemma~\ref{lem:intervals} shows that the number of valid sources in this chain cannot decrease.
Summing over source chains proves that $N_t$ cannot decrease.

For primes $p<q$, the second operation uses the finite chains
\begin{equation}\label{eq:pq-chain}
bp^D,bp^{D-1}q,\ldots,bq^D,
\qquad \gcd(b,pq)=1,\quad D\ge0.
\end{equation}
Again, replace the occupied positions by the same number of initial positions.
Fix a source chain and write $a=bp^{D-e}q^e$, where $0\le e\le D$.
Group the multipliers $j\le t$ according to their unique expressions
\[
j=cp^{d-h}q^h,
\qquad \gcd(c,pq)=1,\quad 0\le h\le d.
\]
For fixed $c,d$ with $cp^d\le t$, the allowed $h$ form an initial interval $0,\ldots,L_{c,d}$, because $p<q$.
The products $aj$ lie at positions $e+h$ in the target chain with parameters $bc,D+d$.
Their membership conditions are thus
\[
e+\{0,\ldots,L_{c,d}\}\subseteq E_{c,d}.
\]
Here $\gcd(bc,pq)=1$, distinct pairs $(c,d)$ give distinct target chains, and $E_{c,d}$ is the occupied set of the corresponding chain.
The simultaneous compression replaces each of these sets by its initial interval.
Lemma~\ref{lem:intervals} again applies.
The group $c=1,d=0$ includes the requirement $a\in S$ itself; before and after compression it restricts $e$ to the source chain, so no additional starts with $e>D$ are counted.
Consequently this operation also cannot decrease any $N_t$.

Each operation preserves cardinality.
Whenever an operation changes the set, it strictly decreases the sum of its elements, since both types of chain are listed in increasing numerical order.
Repeatedly apply any operation that changes the set.
The process must stop because that sum is a positive integer.
Stability under the first operations gives (i), and stability under the second gives (ii).
The inequalities for $N_t$ hold throughout the process.
\end{proof}

\newpage
\section{The upper bound}

For a positive integer $b$, let $\rad(b)$ be the product of its distinct prime divisors, with $\rad(1)=1$.

\begin{lemma}\label{lem:radical}
If $T$ satisfies (i) and (ii) of Lemma~\ref{lem:compression} and $|T|=n\ge3$, then every $b\in T$ satisfies
\[
\log\rad(b)=O(\log n\log\log n),
\]
with an absolute implied constant.
\end{lemma}

\begin{proof}
The case $b=1$ is immediate.
Write the distinct prime divisors of $b$ as $p_{i_1},\ldots,p_{i_m}$.
All their squarefree products belong to $T$, so
\begin{equation}\label{eq:prime-count}
2^m\le n,\qquad m\le\log_2 n.
\end{equation}
The product $p_{i_1}\cdots p_{i_m}$ belongs to $T$.
By successively replacing its factors by smaller primes, every product
\[
p_{j_1}\cdots p_{j_m},\qquad 1\le j_\ell\le i_\ell,
\]
belongs to $T$ as well, including products with repeated prime factors.
There are $\prod_{\ell=1}^m i_\ell$ such ordered choices.
Unique factorization shows that any resulting integer has at most $m!$ preimages, corresponding to permutations of its prime factors.
It follows that
\begin{equation}\label{eq:index-product}
n\ge\frac{\prod_{\ell=1}^m i_\ell}{m!},
\qquad
\sum_{\ell=1}^m\log i_\ell\le\log n+\log(m!).
\end{equation}
Using \eqref{eq:nth-prime}, \eqref{eq:prime-count}, and \eqref{eq:index-product}, we get
\begin{align*}
\log\rad(b)
&=\sum_{\ell=1}^m\log p_{i_\ell}
=O\!\left(m+\sum_{\ell=1}^m\log i_\ell\right)\\
&=O\bigl(\log n+m\log(m+1)\bigr)
=O(\log n\log\log n).
\end{align*}
\end{proof}

Let $T$ be the compression of an arbitrary $n$-element set $S$.
If $p\le r_T(a)$ is prime, then $pa\in T$.
Double counting the pairs $(a,p)$ with $a,pa\in T$, weighted by $\log p$, gives
\begin{align}
\sum_{a\in T}\vartheta(r_T(a))
&\le \sum_{\substack{a\in T,\ p\text{ prime}\\pa\in T}}\log p\notag\\
&=\sum_{b\in T}\sum_{p\mid b}\log p
=\sum_{b\in T}\log\rad(b)
=O(n\log n\log\log n).\label{eq:double-count}
\end{align}
The equality holds because $b/p\in T$ for every prime divisor $p$ of $b\in T$.
By \eqref{eq:primes}, $r=O(1+\vartheta(r))$ for integers $r\ge1$.
Combining \eqref{eq:tails}, Lemma~\ref{lem:compression}, and \eqref{eq:double-count}, we conclude that
\[
\sum_{a\in S}k_S(a)
\le n+\sum_{a\in T}r_T(a)
=O(n\log n\log\log n).
\]

\newpage
\section{A matching construction}

Fix a sufficiently large integer $y\ge4$.
There are $q=\pi(\sqrt y)$ primes at most $\sqrt y$ and $m=\pi(y)-\pi(\sqrt y)$ primes in $(\sqrt y,y]$.
Set
\[
L=\lfloor\log_2 y\rfloor,\qquad H=2qL.
\]
Let $S_y$ consist of the integers
\begin{equation}\label{eq:construction}
a=\prod_{p\le\sqrt y}p^{u_p}\prod_{\sqrt y<p\le y}p^{v_p},
\qquad
0\le u_p\le H,\qquad
\sum_{\sqrt y<p\le y}v_p\le m,
\end{equation}
where all exponents are nonnegative integers.
Unique factorization and stars-and-bars give
\begin{equation}\label{eq:size}
N(y):=|S_y|=(H+1)^q\binom{2m}{m}.
\end{equation}

Call an element good if
\begin{equation}\label{eq:good}
u_p\le H-L\quad(p\le\sqrt y),
\qquad
\sum_{\sqrt y<p\le y}v_p\le m-1.
\end{equation}
If $a$ is good and $1\le j\le y$, then $aj\in S_y$.
Indeed, each small-prime exponent in $j$ is at most $L$.
Also, $j$ has at most one prime factor greater than $\sqrt y$, counting multiplicity, since the product of two such factors exceeds $y$.
The allowances in \eqref{eq:good} therefore suffice for every such $j$, giving $r_{S_y}(a)\ge y$.

By Bernoulli's inequality, the fraction of small-prime exponent choices satisfying \eqref{eq:good} is
\[
\left(\frac{H-L+1}{H+1}\right)^q
=\left(1-\frac{L}{H+1}\right)^q
\ge1-\frac{qL}{H+1}\ge\frac12.
\]
The fraction of large-prime exponent choices satisfying it is
\[
\frac{\binom{2m-1}{m}}{\binom{2m}{m}}=\frac12.
\]
Thus at least a quarter of $S_y$ is good, and
\begin{equation}\label{eq:lower-y}
\sum_{a\in S_y}k_{S_y}(a)\ge\frac{N(y)y}{4}.
\end{equation}

The prime estimates yield $m=\Theta(y/\log y)$.
Moreover $q\le\sqrt y$ and $H=O(\sqrt y\log y)$, so
\[
q\log(H+1)=O(\sqrt y\log y)=o(y/\log y).
\]
Since $\log\binom{2m}{m}=\Theta(m)$, equation~\eqref{eq:size} implies
\begin{equation}\label{eq:asymptotic-size}
\log N(y)=\Theta\!\left(\frac{y}{\log y}\right).
\end{equation}

\newpage
To construct a set of any sufficiently large prescribed size $n$, fix $B>0$ such that $\log N(y)\le By/\log y$ for all large $y$.
Choose a fixed $0<\varepsilon\le1/(2B)$ and put
\[
y=\lfloor\varepsilon\log n\log\log n\rfloor.
\]
For all sufficiently large $n$, we have $y\ge\log n$ and the construction applies, so
\[
\log N(y)\le\frac{By}{\log y}
\le B\varepsilon\log n\le\tfrac12\log n.
\]
Thus $N(y)\le\sqrt n\le n/2$.
Put $s=\lfloor n/N(y)\rfloor$, and take $s$ disjoint dilates of $S_y$.
For example, if $Q>\max S_y$, the sets $Q^iS_y$ for $0\le i<s$ are disjoint.
Their union contains $sN(y)>n-N(y)\ge n/2$ elements; pad it with arbitrary new positive integers to reach cardinality $n$.
At least $sN(y)/4\ge n/8$ elements retain all their first $y$ multiples, since dilation preserves these relations and adding elements cannot destroy them.
Consequently
\[
\sum_{a\in S}k_S(a)\ge\frac{ny}{8}
=\Omega(n\log n\log\log n).
\]
This completes the proof of Theorem~\ref{thm:main}.

\appendix
\section{The prime estimates used above}

For completeness, the bounds \eqref{eq:primes} follow from central binomial coefficients.
Every prime in $(n,2n]$ divides $\binom{2n}{n}$, hence
\[
\vartheta(2n)-\vartheta(n)\le\log\binom{2n}{n}\le2n\log2.
\]
Summing over dyadic intervals and using monotonicity gives $\vartheta(x)=O(x)$.

Define $\psi(x)=\sum_{p^h\le x}\log p$.
For each prime $p$, the exponent of $p$ in $\binom{2n}{n}$ is
\[
\sum_{h\ge1}\left(\left\lfloor\frac{2n}{p^h}\right\rfloor
-2\left\lfloor\frac{n}{p^h}\right\rfloor\right).
\]
Each summand is zero or one, and vanishes when $p^h>2n$.
It follows that
\[
\psi(2n)\ge\log\binom{2n}{n}\ge2n\log2-\log(2n+1),
\]
where the last inequality uses the fact that the central binomial coefficient is the largest of the $2n+1$ coefficients whose sum is $2^{2n}$.
By monotonicity, $\psi(x)\ge c_0x$ for some $c_0>0$ and all sufficiently large $x$.
The terms with $h\ge2$ involve only primes at most $\sqrt x$, and each such prime contributes at most $\log x$ to their total.
Thus
\[
0\le\psi(x)-\vartheta(x)\le\sqrt x\log x=o(x),
\]
which proves $\vartheta(x)\asymp x$.
Finally,
\[
\frac{\vartheta(x)}{\log x}\le\pi(x)
\le\sqrt x+\frac{2\vartheta(x)}{\log x},
\]
giving $\pi(x)\asymp x/\log x$.
In particular, $\pi(x)\ge\sqrt x$ for all sufficiently large $x$.
Taking $x=p_j$ gives $j\ge\sqrt{p_j}$, hence $\log p_j\le2\log j$ for all sufficiently large $j$.
Enlarging the constant to cover the remaining indices proves \eqref{eq:nth-prime}.

\end{document}
